%-------------------------------------------------------------------
\section{Model}\label{sec:model}

\paragraph{Code and noise.}
We use dimensionless quadratures with $[\hat q,\hat p]=i$ and the square GKP
code. Its stabilizer $S_q=e^{2i\sqrt\pi\hat q}$ fixes the lattice spacing
$L=\sqrt\pi$. A shift by $L$ is the logical $\bar X$ and a shift by $2L$ is a
stabilizer. As in the proposal, we model only the $q$ quadrature and simulate
classical shifts rather than the oscillator state. For pure displacement
noise this reproduces the ideal-code quantum dynamics; the simulator itself is
validated separately in \texttt{docs/sim\_validation.pdf}.
In round $t=1,\dots,50$ of a trial,
\begin{align*}
  x &\leftarrow x+\varepsilon_t, &
  s_t &= \operatorname{wrap}(x+\nu_t)\in[-L/2,L/2), &
  x &\leftarrow x-c_t,
\end{align*}
where $c_t$ is the decoder's correction. The displacement
$\varepsilon_t$ follows a stationary AR(1) process,
$\varepsilon_t=\rho\varepsilon_{t-1}+\eta_t$ with
$\operatorname{Var}\varepsilon_t=\sigma_q^2$, and $\nu_t\sim\mathcal N(0,\sigma_m^2)$
is i.i.d.\ measurement noise. The nominal point is $\sigma_q=0.2L$,
$\sigma_m=0.05L$, $\rho=0.9$. After round $t$ the qubit is in a logical error
when $\operatorname{round}(x/L)$ is odd. The \emph{logical error rate}
$P_L(n)$ is the fraction of trials in error after round $n$.

\paragraph{The decoding problem.}
Every decoder in this study applies $c_t=s_t+j_tL$ with $j_t\in\{0,1\}$; the
nearest-lattice (NL) decoder always plays $j_t=0$. Under this restriction the
residual after each round is an integer multiple of $L$ minus $\nu_t$, so
\begin{equation}
  s_t=\operatorname{wrap}(z_t),\qquad z_t=\varepsilon_t+\nu_t-\nu_{t-1},
  \label{eq:z}
\end{equation}
\emph{whatever the decoder did earlier}. Let
$y_t=\operatorname{round}(z_t/L)\bmod 2$ be the parity of the NL decoder's
jump in round $t$, and $K_t=\sum_{\tau\le t}y_\tau\bmod 2$ the NL decoder's
logical state. Then any decoder's logical state after round $t$ is
\begin{equation}
  K_t\oplus J_t,\qquad J_t=\sum_{\tau\le t}j_\tau\bmod 2 .
  \label{eq:state}
\end{equation}
Decoding therefore amounts to estimating the parity $K_t$ from the wrapped
history $s_1,\dots,s_t$. The best decoder plays $J_t=\hat K_t$, the MAP
estimate given all syndromes so far. Because later syndromes do not depend on
the $j$'s, this greedy rule minimises $P_L(t)$ for every $t$ at once, and it
lets a decoder \emph{revise} an earlier decision once later syndromes show it
was wrong. A test in \texttt{tests/test\_core.py} checks
\eqref{eq:z}--\eqref{eq:state} against the simulator for an arbitrary decision
sequence.

Two consequences shape the study:
\begin{itemize}
  \item Every decoder sees the same syndromes on the same noise, so all
    comparisons are paired on common random numbers.
  \item Supervised labels for the LSTM ($y_t$ and $K_t$ from an NL run of the
    simulator) do not depend on the LSTM's own decisions, so training data
    and closed-loop testing are not mismatched.
\end{itemize}
Continuous fine corrections (estimating $\nu_t$) are deliberately excluded.

%-------------------------------------------------------------------
\section{Decoders}\label{sec:decoders}

\paragraph{Nearest lattice (NL).} $c_t=s_t$, i.e.\ $j_t=0$. Uses only the
current syndrome.

\paragraph{Bayesian grid filter.}
The hidden state is $(X_t \bmod 2L,\ \varepsilon_t)$, where
$X_t=\sum_{\tau\le t}\varepsilon_\tau$ is the cumulative uncorrected
displacement. The decoder knows its cumulative correction $C_{t-1}$, and
$s_t=\operatorname{wrap}(X_t-C_{t-1}+\nu_t)$. This is an exact hidden Markov
model of the simulator. On a grid with spacing $h$, each round performs:
\begin{enumerate}
  \item A bin-integrated Gaussian AR(1) prediction of $\varepsilon_t$ (one
    matrix product). The innovation variance is floored at $h^2/12$.
  \item A shear $X_t=X_{t-1}+\varepsilon_t$ (an integer index shift per
    $\varepsilon$ row).
  \item Multiplication by the periodic likelihood
    $\sum_m\phi_{\sigma}\bigl(s_t+C_{t-1}-X+mL\bigr)$, evaluated analytically
    with $\sigma^2=\sigma_m^2+h^2/12$.
  \item Normalisation.
\end{enumerate}
The correction class $j_t$ is the one whose half of the $2L$ circle around
$s_t+C_{t-1}+j_tL$ carries more posterior mass. This is the cumulative MAP
rule above. The filter is batched over trials on the GPU, and its grid
convergence is shown in Table~\ref{tab:e0}.

We use the filter with several assumed models:
\begin{itemize}
  \item \emph{Bayes (matched)} uses the true parameters. It is the
    syndrome-only optimum up to grid error.
  \item \emph{Bayes (nominal)} always assumes the nominal parameters, which
    makes it the counterpart of an LSTM trained only at the nominal point.
  \item \emph{Bayes-broad} is a Bayesian model average over 8 filters with
    $\rho'\in\{0,0.6,0.85,0.95\}$ and $\sigma_q'\in\{0.15,0.25\}L$, weighted
    by their marginal likelihood. It is the counterpart of the broadly
    trained LSTM.
  \item \emph{Bayes (regime-aware)} adds the hidden correlation regime to the
    state, for the switching-noise experiment.
  \item \emph{Bayes (per-round)} is an ablation that, after each decision,
    discards the posterior mass of the rejected class, so it never revises a
    past decision.
\end{itemize}

\paragraph{LSTM.}
A two-layer LSTM (hidden size 128) reads $[s_t/L,\ \sin(2\pi s_t/L),\
\cos(2\pi s_t/L)]$ and has five sigmoid heads. Head ``$k$'' predicts the
window parity $W^{(k)}_t=y_{t-k}\oplus\dots\oplus y_t$ for
$k\in\{0,1,2,4\}$, and head ``full'' predicts $K_t$. A head's prediction is
turned into actions through $J_t=\hat W^{(k)}_t\oplus J_{t-k-1}$, with
$J_t=\hat K_t$ for the full head. The $k=0$ head is the per-round
(no-revision) ablation.

Training uses 200{,}000 sequences of 50 rounds, summed binary cross-entropy
over all heads, AdamW with a cosine schedule, at most 40 epochs, and early
stopping on validation loss. The head used at test time is the one with the
lowest validation time-averaged logical error, and test data is never used
for any choice. The training sets are:
\begin{itemize}
  \item \emph{nominal}: the nominal point only, 3 seeds;
  \item \emph{broad}: per-sequence $\rho\sim U(0,0.97)$,
    $\sigma_q\sim U(0.12,0.28)L$, $\sigma_m\sim U(0.02,0.08)L$, 3 seeds;
  \item \emph{matched}: two single-point models at $\rho=0.5$ and at
    $\sigma_q=0.26L$, used to measure the generalisation gap;
  \item \emph{switching}: the switching-noise model of E6.
\end{itemize}
Labels come from the true errors through an NL run (proposal, Sec.~5).
At test time the LSTM sees only syndromes.

%-------------------------------------------------------------------
\section{Evaluation protocol}\label{sec:protocol}
\begin{itemize}
  \item \textbf{Shared test noise.} At each test point all decoders run on
    the same test noise. Random streams are derived from a
    \texttt{SeedSequence} keyed by (split, experiment, point, chunk), so
    train, validation, development and test streams are disjoint; this is
    checked by a test.
  \item \textbf{Test-set size.} The size is set once per point: 20{,}000-trial
    chunks are added until matched Bayes shows at least 50 errors at round 50,
    or a cap is reached. The caps are 100{,}000 trials for E1, 60{,}000 for E6,
    40{,}000 for E2 and E5, and 20{,}000 for E3 and E4.
  \item \textbf{Reported quantities.} We report $P_L(10)$ and $P_L(50)$ with
    Wilson 95\% intervals, plus the time-averaged logical error
    $\frac1{50}\sum_{n=1}^{50}P_L(n)$. When no errors are observed we report
    the 95\% upper bound $3/N$.
  \item \textbf{Paired differences.} Differences between decoders are means
    of per-trial differences, with their standard errors.
  \item \textbf{LSTM seeds.} LSTM results are averaged over seeds, and the
    spread across seeds is given separately from the test-noise error.
  \item \textbf{In or out of distribution.} Points outside the broad training
    box are marked (hollow markers in figures, italics in tables).
\end{itemize}
